\section{Tartışma}\label{sec:tartisma}

\paragraph{Ucuz bir vekil neden önde kalır?}
Gözlenen performans profili, birbirini destekleyen üç temel mekanizmayla tutarlıdır. Birincisi
hedef uyumu hipotezidir. Double-bridge vuruşunun getirisi, kırılan kenarın yerine daha kısa
bağlantılar kurulabilmesine bağlıdır; uzun bir kenar, özellikle uç noktalarının komşu ölçeğine
göre uzunsa, tam da bu kazancı vaat eder. Üyelik olasılığı ise ``bu kenar optimal turda mı?''
sorusunu yanıtlar ve havuzu, düzeltildiğinde az kazandıran hatalarla doldurabilir. Aşama~2 bu
hipotezi desteklerken holdout2'de GNN ile kNN ve kazanç-MLP ile üyelik-MLP ayrışmamıştır
(ikisinde de $p_{\mathrm{Holm}}=1{,}0$; Bölüm~\ref{sec:bulgular-asama2}).
İkincisi maliyettir: skor sabit zamanlıdır ve bütçenin neredeyse tamamı vuruşa ve yerel aramaya
kalır. Üçüncüsü tazeliktir: ucuz skor her tur değişiminde yeniden hesaplanabildiği için kesimler
güncel tura düşer. Bu yaklaşım, literatürdeki köklü sezgisellerle kuramsal bir süreklilik
içindedir. GLS ceza hedefini kenar maliyetine göre seçerek uzun kenarları önceliklendirir
\citep{voudouris_1999_gls}; rotalamada istenmeyen kenar özniteliklerinden türetilen bilginin
pertürbasyonu yönlendirebildiği de gösterilmiştir \citep{arnold_2019_kgls}. Desen, basit tabanların
öğrenilmiş rehberlerle rekabet ettiği önceki gözlemlerle aynı yöndedir
\citep{xia_2024_rethinking, bother_2022_wrong}.
% kaynak: E605, E606
Eşit vuruş oturumu bu üç mekanizmayı birbirinden ayırır (Bölüm~\ref{sec:bulgular-hakem}). Vuruş
sayısı eşitlendiğinde öğrenilmiş skorlar vekilden kötü değildir; kazanç-MLP ve NeuroLKH biraz öndedir.
Eşit süredeki sıralama bu yüzden ağırlıkla maliyet ve tazelikten doğar; hedef uyumu tek başına
vekilin üstünlüğünü açıklamaz.
% kaynak: E1128-E1130, E1137, E1139

\paragraph{Hangi eğitimsiz sinyal?}
Eğitimsiz olmak tek başına yetmez. kNN sırası holdout2'de rastgeleden anlamlı ayrışmadı ve karışık
problemlerde gerisinde kaldı (betimsel); fazla-uzunluk ise kNN'yi 54 problemin 43'ünde geçti
(keşifsel).
kNN bir kenarın komşuluk \emph{sırasını}, fazla-uzunluk ise kenarın \emph{büyüklüğünü} ölçer;
her iki kol geliştirme setinde kendi en iyi havuz kuralıyla yapılandırıldığından, bu karşılaştırma
pratik uygulama düzeyindeki uçtan uca strateji başarımını yansıtır.
Keşifsel ham uzunluk kolu, yerel ölçek düzeltmesinin ablasyonudur ve ``önce en basit vekile karşı sına'' ilkesini kendi vekilimize uygular;
büyüklük bilgisinin temel belirleyici olduğunu düşündürür: yerel ölçek düzeltmesi olmadan en uzun
kenarları hedeflemek de rastgeleye karşı benzer kazanç sağladı ve vekilden ortalamada ayrışmadı
(fark $+0{,}009$ puan, \%95 GA $[-0{,}032;\,+0{,}053]$).
Eşit vuruş oturumu bu tabloyu inceltir: vuruş başına en iyi sinyal kNN sırasıdır (vekilden
$0{,}088$ puan iyi, 37/54). kNN'nin eşit sürede geride kalması, sinyalinden çok vuruş hızından
kaynaklanır (aynı oturumda 867'ye karşı 1.196 vuruş/sn). İki eğitimsiz sinyal arasındaki seçim bu
nedenle bir nitelik--maliyet takasıdır.
Uygulayıcı için sonuç sadedir: bu motorda ve 10~sn'de, komşu listesi dahi gerektirmeyen ham kenar
uzunlukları güçlü bir başlangıç noktasıdır. Yerel ölçek düzeltmesi ortalamada ek kazanç üretmemiş
olup, kümeli yapılardaki olası avantajı bağımsız bir doğrulayıcı hipotez olarak değerlendirilebilir.
% kaynak: E604, E750, E751, E965, E972, E973, E976, E1131, E1137, E1145

\paragraph{Hedefleme nerede işe yarar?}
Kazanç kümeli ve gerçek problemlerde toplanır, uniform problemlerde küçük kalır: hedefleme bir
``kötü kenar'' sinyaline dayanır ve bu sinyal ancak kenar uzunlukları heterojen olduğunda bilgi
taşır. Katman etiketinden bağımsız kümelenme göstergesi de bunu destekler: kazanç, en yakın komşu
uzaklığının değişim katsayısıyla birlikte büyür ($\rho=-0{,}56$).
Yalnız uniform örneklerden oluşan bir kıyas kümesi bu yüzden hedeflemenin değerini
olduğundan küçük gösterir.
% kaynak: E1038

\paragraph{Yayımlanmış bir modele ve daha güçlü bir yerel aramaya karşı.}
NeuroLKH'nin sinyali, süresi sayılmadığında problemlerin çoğunda vekilden küçük bir farkla iyidir
(medyan $0{,}066$ puan). Öğrenilmiş bir modelin taşıdığı ek bilgi gerçektir, ama küçüktür; tek
seferlik ve kısa bütçeli bir çözümde ağın süresi bu kazancı siler. Model çıktısının birçok koşuya
yayılabildiği durumlarda, örneğin aynı problemin tekrar tekrar çözülmesinde, NeuroLKH gibi bir
model küçük bir kazanç sağlayabilir. Tek seferlik kullanımda eğitimsiz vekil daha basittir ve
ondan istatistiksel olarak ayrışmaz.
LKH-3'te ise hedefleme ölçülebilir bir kazanç sağlamadı. Kayıtta öngörülen olası açıklama şudur:
LK'nin derin hamleleri uzun kenarların çoğunu zaten kendisi onarır. Vuruş hedeflemenin değeri bu
yüzden yerel aramanın gücüne bağlıdır; bulgular hızlı ve görece zayıf yerel aramalı ILS için
geçerlidir.
% kaynak: E1106, E1107, E1172, E1180

\paragraph{Çevrimiçi öğrenme ve ödül tasarımı.}
Çevrimiçi model başarısız bir öğrenici değildir; kendisine verilen amacı yükseltir. Sorun amaçtadır:
vuruşların çoğu turu iyileştirmediğinde vuruş başına ortalama ödül, yerel aramanın onardığı nötr
vuruşu ödüllendirir (Bölüm~\ref{sec:bulgular-mekanizma}). Uyarlamalı operatör seçimi literatürü, kredi atamasının seçim başarımını etkilediğini \citep{fialho_2010_aos} ve nadir ama büyük iyileşmeleri öne çıkaran uç-değer kredisinin ortalama krediye göre avantajlı olabileceğini \citep{fialho_2008_extreme} öne sürer. Geliştirme setindeki ölçümümüz (12 problem, keşifsel) bunun bir örneğidir. Hakem odaklı blok bunu
holdout2'de doğrudan sınadı: kredi $\max(0,r)$ ya da kabul göstergesi olduğunda zarar 54 problemin
54'ünde kaybolur ve model sabit ön-bilgi düzeyine döner. Tasarım dersi şudur: seyrek iyileşmeli
aramalarda kredi, işaretli ortalama maliyet değişimi yerine kabul edilen iyileşmelere dayanmalı ve
öğrenicinin kendi ödülü son kaliteyle birlikte izlenmelidir. Zarar öğrenmenin kendisinden değil
kredi tanımından gelir; doğru krediyle de bu bütçede sabit modele bir kazanç eklenmedi.
% kaynak: E1148-E1152

\paragraph{Uygulayıcı için öneriler.}
\begin{itemize}
\item Vuruşları en uzun tur kenarlarına yönlendirin (ham ya da uç noktaların yakın-komşu ölçeğine
göre düzeltilmiş uzunluk); skor sabit zamanlıdır, eğitim ve çözücü etiketi gerektirmez. Bu tür
motorlarda hedefli global double-bridge, rastgele segment vuruşundan belirgin biçimde iyidir.
% kaynak: E1047
\item Havuzu tur kenarlarının \%4--6'sı kadar, hedefleme olasılığını $0{,}5$--$0{,}75$ aralığında
tutun; yalnız küçük havuz ile çok yüksek olasılığın birleşiminden kaçının. Bu plato holdout2'de de
görüldü.
% kaynak: E750-E752, E1091
\item En büyük kazancı kümeli ve gerçek (VLSI ve ulusal) örneklerde bekleyin; GNN'ye karşı fark
kısa bütçelerde en büyüktür.
\item Öğrenilmiş bir hedefleme skorunu bu vekile karşı, model süresi bütçeye dahil ve eşit
duvar-saatinde sınayın; sinyalin niteliğini ayrıca eşit vuruş sayısında ölçün. AUC gibi kestirim
ölçütleri arama faydasının yerine geçmez.
\item Arama içinde öğrenme kullanacaksanız krediyi kabul edilen iyileşmelere dayandırın
($\max(0,r)$ ya da kabul göstergesi), işaretli ortalama ödülden kaçının ve ağırlıkların iyi bir
statik ön-bilgiden ne kadar saptığını izleyin.
\item LK/LKH düzeyinde bir yerel arama kullanıyorsanız vuruş hedeflemeden kazanç beklemeyin; LKH-3'te
etkisi ölçülemedi.
\end{itemize}

\subsection{Kapsam ve geçerlilik sınırları}\label{sec:tartisma-sinir}

Bu çalışmanın doğrulayıcı bulguları, 2-opt/Or-opt tabanlı yüksek frekanslı ILS motoru, 10~sn
duvar-saati bütçesi ve $n=1.689$--$9.882$ ölçek aralığı için tesis edilmiştir.
Lin--Kernighan (LK) düzeyindeki yerel aramada hedeflemenin katkısı keşifsel olarak sınandı:
LKH-3'ün ILS kipinde hedefli vuruş rastgele vuruştan ayrışmadı. LKH'de yalnız tek bir yama, tek
bir bütçe ve ILS kipi denendi; bulgular 2-opt/Or-opt tabanlı hızlı ILS için geçerlidir.
Benzer biçimde segment vuruşu bu motorda sabit $W=50$ penceresiyle değerlendirilmiş olup, LK
tabanlı çözücülerdeki davranışı kendi mimari bağlamında ele alınmalıdır.
Keşifsel 60~sn bloğu (5 tohum), vekilin GNN'ye karşı avantajının daha uzun bütçelerde kademeli
olarak daralmakla birlikte korunduğunu göstermektedir ($-0{,}078$ puan, \%95 GA üst sınırı
$-0{,}008$); dakika mertebesini aşan bütçelerdeki asimptotik denge gelecek çalışmalarda
incelenebilir.
Donanım yükü dalgalanmalarının etkisini dışlamak amacıyla, bütün strateji kıyasları aynı oturum ve
aynı işçi sürecinde eşleşik olarak yürütülmüştür (ana oturumda rastgele global kol 3.252, yerel
vuruş oturumunda 2.225, ham uzunluk oturumunda 3.269 vuruş/sn; vekil sırasıyla 2.288 ve
2.773 vuruş/sn).
Kol sırası ana koşuda tohumlar arasında döndürüldü; ilk sıradaki kolun yaklaşık \%18 daha yavaş
vurduğu sonradan ölçüldü. Tam rastgele sırayla ve tek işçiyle koşulan bloklarda da vekil rastgeleden
düşük gap verdi.
% kaynak: E597, E598, E697, E705, E986, E987, E1057, E1078, E1113, E1164, E1181

Karşılaştırılan GNN modeli, $n\le400$ sentetik örneklerde eğitilmiş hafif bir mimari olarak büyük
test örneklerinde ($n\le9.882$) boyut genellemesi senaryosunu temsil eder; doğrudan hedef boyutta
eğitilmiş daha yüksek kapasiteli ağların hız--kalite takası ayrı bir değerlendirme konusudur.
Yayımlanmış ve daha büyük veriyle eğitilmiş NeuroLKH ile yapılan keşifsel karşılaştırma bu sınırı
kısmen kapatır: ağın süresi sayıldığında vekil ondan kötü değildir, sayılmadığında NeuroLKH
problemlerin çoğunda küçük bir farkla öndedir.
V4 testinde kazanç-MLP varsayılan havuz kuralıyla, vekil ise geliştirme setinde belirlenen kuralla
koşulmuştur; geliştirme setinde vekilin seçilen ve varsayılan kuralları arasında anlamlı bir fark
bulunmaması ($p=0{,}36$) bu karşılaştırmanın kararlılığını destekler. Kazanç-MLP aynı seçim
kuralıyla ayarlandığında da vekil önde kaldı ($-0{,}070$ puan).
Hedeflemenin etkinliği problem yapısına bağlıdır: kümeli ve gerçek problemlerde belirgin kazançlar
sağlanırken, kenar uzunluklarının homojen olduğu uniform dağılımlarda kazanç beklendiği üzere
sınırlı kalır ve tekil örneklerde (örn.\ h2c13) öğrenilmiş skorlar öne çıkabilir.
Son olarak, ön-kayıt protokolü (OpenTimestamps, Bitcoin blok 976258) uyarınca V1--V4 testleri
doğrulayıcı niteliktedir; katman analizleri, ikincil ikili karşılaştırmalar, ham uzunluk kolu ve
sağlamlık blokları ise bu ana bulguları derinleştiren ve sonraki ön-kayıtlı çalışmalar için
hipotez üreten keşifsel kanıtlar olarak konumlandırılmıştır.
% kaynak: E585, E751, E1107, E1110
